\documentclass[../../main.tex]{subfiles}

\begin{document}

\section{The density matrix, entanglement entropy and their connection to singular value decomposition (SVD)}
\label{sec:density-matrix-entanglement-svd}

\subsection{Main goal}

Prove that, for a bipartite pure state, the von Neumann entanglement entropy is
\begin{equation}
\boxed{\,S=-\Tr\bigl(\rho_A\log\rho_A\bigr)
=-\sum_i\lambda_i\log\lambda_i
=-\sum_i s_i^2\log s_i^2\,},
\end{equation}
where $\lambda_i$ are the eigenvalues of the reduced density matrix $\rho_A$
and $s_i$ are the Schmidt values, with $\lambda_i=s_i^2$.
We use the convention $0\log 0=0$.

\subsection{The density operator}

Let $\{p_i,\ket{\psi_i}\}$ be an ensemble of pure states, describing a system that is in one
of the states $\ket{\psi_i}$ with probability $p_i$, i.e.\ the system is not fully known. Define
the density operator as
\begin{equation}
\rho=\sum_i p_i\,\dyad{\psi_i}{\psi_i},\qquad p_i\in\mathbb R\ \, \&\ \sum_i p_i=1,
\end{equation}
which is a positive semi-definite operator of trace one (see for example Nielsen \& Chuang).

\subsection{Observables on a composite system}

Take a system $A$ with some observable $M$, and a system $B$, giving the composite system
$AB$ (see Nielsen \& Chuang Box 2.6). Assume we can measure $M$ on system $A$; then a corresponding measurement
on $AB$ is $\tilde M=M\otimes\mathds{1}_B$.

\begin{proof}
Take $AB$ to be prepared in $\ket{m}\ket{\psi}$ with $M\ket{m}=m\ket{m}$. Applying the
measurement to $\ket{m}\ket{\psi}$ should yield $m$ with probability one, i.e.\ the
corresponding projector for $\tilde M$ is $P_m\otimes\mathds{1}_B$. By the spectral theorem,
$\tilde M=\sum_m m\,P_m\otimes\mathds{1}_B=M\otimes\mathds{1}_B$ for a general state of $AB$.
\end{proof}

\subsection{The reduced density matrix as a partial trace}

Measurement averages must be the same whether computed via the full $\rho^{AB}$ or via some
$\rho^{A}$ associated to system $A$, i.e.
\begin{equation}
\Trace\bigl(M\rho^A\bigr)\overset{!}{=}\Trace\bigl(\tilde M\rho^{AB}\bigr)
=\Trace\bigl((M\otimes\mathds{1}_B)\rho^{AB}\bigr).
\end{equation}
Take $\rho^A=\Trace_B(\rho^{AB})$ where
\begin{equation}
\Trace_B\bigl(\dyad{a_1}{a_2}\otimes\dyad{b_1}{b_2}\bigr)
\equiv\dyad{a_1}{a_2}\,\Trace\bigl(\dyad{b_1}{b_2}\bigr)
\qquad\&\qquad\Trace_B(\cdot)\ \text{linear.}
\end{equation}
Write $\rho^{AB}$ by using the eigenbasis of $M$ for system $A$ with $M\ket{i}=m_i\ket{i}$ and
$\braket{i}{j}=\delta_{ij}$:
\begin{align}
\rho^{AB}&=\sum_\alpha p^\alpha\dyad{\psi}{\psi}^\alpha
=\sum_\alpha p^\alpha\sum_{ijk\ell}c_{ij}^\alpha(c_{k\ell}^\alpha)^*\,
\ket{i}_A\ket{j}_B\bra{\ell}_B\bra{k}_A,\\
\rho^{A}&=\sum_\alpha\sum_{ijk\ell}p^\alpha c_{ij}^\alpha(c_{k\ell}^\alpha)^*\,
\dyad{i}{k}_A\underbrace{\Trace\bigl(\dyad{j}{\ell}_B\bigr)}_{=\delta_{j\ell}}
=\sum_{\alpha,ijk}p^\alpha c_{ij}^\alpha(c_{kj}^\alpha)^*\,\dyad{i}{k}_A.
\end{align}
Therefore
\begin{align}
\Trace\bigl(M\rho^A\bigr)
&=\sum_{\alpha,ijk}p^\alpha c_{ij}^\alpha(c_{kj}^\alpha)^*\,
\underbrace{\Trace\bigl(m_i\dyad{i}{k}_A\bigr)}_{=m_i\delta_{ik}}
=\sum_{\alpha,ij}p^\alpha m_i\,c_{ij}^\alpha(c_{ij}^\alpha)^*,
\end{align}
and similarly $\Trace\bigl((M\otimes\mathds{1}_B)\rho^{AB}\bigr)
=\sum_{\alpha,ij}p^\alpha m_i\,c_{ij}^\alpha(c_{ij}^\alpha)^*$, which proves the claim. One can
also easily show that this is the only map that conserves consistency (see Nielsen \& Chuang Box 2.6).

\subsection{Reduced density matrix of a pure state}

Now examine the reduced density matrix for a pure state,
\begin{align}
\rho_{AB}&=\dyad{\psi}{\psi}
=\sum_{ijk\ell}c_{ij}(c_{k\ell})^*\,\ket{i}_A\ket{j}_B\bra{\ell}_B\bra{k}_A,\\
\rho_A&=\Trace_B(\rho_{AB})=\sum_{ijk}c_{ij}(c_{kj})^*\,\dyad{i}{k}_A, \label{eq:reduced_density_matrix_as_operator}
\end{align}
such that the matrix representation is
\begin{equation}
\boxed{\,\rho_A=\psi\psi^\dagger\,}, \label{eq:density_matrix_as_wavefunction_coefficients}
\end{equation}
using that the $c_{ij}$ parametrize the matrix $\psi\in\mathbb C^{d_A\times d_B}$, i.e.\
$\ket{\psi}=\sum_{ij}c_{ij}\ket{i}_A\ket{j}_B$ with $\sum_{ij} \vert c_{ij} \vert^2 = 1$.


Furthermore, \cref{eq:reduced_density_matrix_as_operator} results in
\begin{align}
      \Trace \rho_A &= \Trace \left[ \sum_{ijk} c_{ij} (c_{kj})^* \dyad{i}{k}_A \right] \notag \\
      &= \sum_{ijkl} c_{ij} (c_{kj})^* \delta_{il} \delta_{kl} \notag \\
      &= \sum_{jl} c_{lj} (c_{lj})^* = 1 , \notag
\end{align}
which means that that the eigenvalues of the reduced density matrix sum to one. This can be seen using that any hermitian matrix can be diagonalized and the cylic property of the trace
\begin{equation}
      \Trace \rho_A = \Trace [O D O^\dagger] = \Trace D = \sum_i \lambda_i = 1. \label{eq:eigenvalues_reduced_density_matrix_sum_to_one}
\end{equation}


\subsection{Connection to the SVD}

Now use the SVD on $\psi$:
\begin{equation}
\psi=USV^\dagger\qquad\text{with }U\in\mathbb C^{d_A\times d_A},\ V\in\mathbb C^{d_B\times d_B}\ \text{unitary.}
\end{equation}
Then, using $V^\dagger V=\mathds 1$,
\begin{equation}
\rho_A=USV^\dagger VS^\dagger U^\dagger
=USS^\dagger U^\dagger,
\qquad
SS^\dagger=\diag(s_1^2,\dots,s_r^2,0,\dots,0),
\end{equation}
and write $S^2 = \diag(s_1^2,\dots,s_r^2)$.
Furthermore, since $\rho_A$ is hermitian we can diagonalize it with
$\rho_A=ODO^\dagger=US^2 U^\dagger$ and $D=\diag(\lambda_1,\dots)$, so that
\begin{equation}
\boxed{\,\lambda_i=s_i^2\,},
\end{equation}
i.e.\ the eigenvalues are given by the squared Schmidt values and hence non-negative. Thus, we can conclude together with \cref{eq:eigenvalues_reduced_density_matrix_sum_to_one} that 
\begin{equation}
      \boxed{\, 0 \leq \lambda_i \leq 1 \,} . \label{eq:range_of_eigenvalues_of_reduced_density_matrix}
\end{equation}

\subsection{Functions of the density matrix via the spectral theorem}

Since $\rho_A$ is hermitian, we can use the spectral theorem to write it as
$\rho_A=\sum_i\lambda_i P_i$ where $\lambda_i$ are its eigenvalues and the eigenvectors
$\ket{u_i}$ give $P_i=\dyad{u_i}{u_i}$, since
\begin{equation}
\rho_A=UDU^\dagger=\bigl(\ket{u_1},\dots,\ket{u_n}\bigr)
\begin{pmatrix}\lambda_1& &\\ &\ddots&\\ & &\lambda_n\end{pmatrix}
\begin{pmatrix}\bra{u_1}\\\vdots\\\bra{u_n}\end{pmatrix},
\qquad\braket{u_i}{u_j}=\delta_{ij}.
\end{equation}


For a function $f$ defined on the spectrum of $\rho_A$, the spectral
functional calculus\footnote{See for example Nicholas, J. H. (2008). \emph{Functions of matrices: theory and computation.}} gives
\begin{equation}
f(\rho_A)\equiv\sum_i f(\lambda_i)\dyad{u_i}{u_i}
=U\diag\bigl(f(\lambda_i)\bigr)U^\dagger.
\end{equation}

\subsection{The von Neumann entanglement entropy}

Therefore, for the von Neumann entanglement entropy,
\begin{align}
S(\rho_A)&\equiv-\Tr\bigl(\rho_A\log\rho_A\bigr) \\
&=-\Tr\bigl(\rho_A\,U\log(\diag(\lambda_i))\,U^\dagger\bigr)\notag\\
&=-\Tr\bigl(U\diag(\lambda_i)\underbrace{U^\dagger U}_{=\mathds 1}\log(\diag(\lambda_i))U^\dagger\bigr)\notag\\
&=-\sum_i\lambda_i\log\lambda_i
=-\sum_i s_i^2\log s_i^2.
\end{align}
$S$ is zero for two unentangled systems and non-zero for entangled systems.

\subsection{The Rényi entanglement entropies}

The Rényi entropy of order $n$ is defined as
\begin{align}
      S_n(\rho_A) &\equiv \frac{1}{1-n} \ln (\Trace[\rho_A^n]) \\
      &=  \frac{1}{1-n} \ln \left( \sum_{i=1}^{d_A} \lambda_i^n \right),
\end{align}
using again, that any hermitian matrix is diagonalizable and the cylic property of the trace. 
Employing l'Hôpital's rule yields for order one that
\begin{align}
      S_1(\rho_A) &= \lim_{n \rightarrow 1} \frac{\ln (\Trace[\rho_A^n])}{1-n} 
      \overset{\mathrm{H}}{=} \lim_{n \rightarrow 1} \frac{\left( \ln \sum_{i=1}^{d_A} \lambda_i^n \right)'}{(1-n)'} \notag \\
      &= - \lim_{n \rightarrow 1} \frac{\sum_{i=1}^{d_A} \lambda_i^n \ln \lambda_i }{\sum_{i=1}^{d_A} \lambda_i^n} 
      = - \frac{\sum_{i=1}^{d_A} \lambda_i \ln \lambda_i }{\sum_{i=1}^{d_A} \lambda_i} \notag \\
      &= - \sum_{i=1}^{d_A} \lambda_i \ln \lambda_i .
\end{align}
That means, that the Rényi entropy of order one is reduces to the von Neumann entropy.


\subsection{Example}

In our case, for example, $L=2$, $N_f=1$, $N_\mathrm{ph}=1$:
\begin{align}
& \{\ket{0},\ket{1}\}\ \&\ \{\ket{01},\ket{10}\}
\quad\text{as bases of }\mathcal H_b\ \&\ \mathcal H_f,\\
&\Rightarrow\{\ket{0, 01};\ \ket{0, 10};\ \ket{1, 01};\ \ket{1, 10}\}
\quad\text{as basis of }\mathcal H_b\otimes\mathcal H_f,\\
&\Rightarrow\ket{\psi}=\sum_{i=1}^{2}\sum_{j=1}^{2}\psi_{ij}\ket{i}_b\ket{j}_f
\quad\text{with }(\psi)_{ij}=\psi_{ij}\ \&\ \psi\in\mathbb C^{2\times2},\\
&\Rightarrow S(\rho_A)=-\sum_\ell s_\ell^2\log s_\ell^2.
\end{align}
Note that since $\Trace\rho^{AB}=\Trace\rho^A=\Trace\rho^B=1$ we must have
$\sum_\ell\lambda_\ell=\sum_\ell s_\ell^2=1$.

\subsection{References}
\begin{itemize}
\item M.\ Nielsen \& I.\ Chuang, \emph{Quantum Computation and Quantum Information}
      (density operator, partial trace; Box 2.6).
\item U.\ Schollw\"ock, \emph{The density-matrix renormalization group in the age of matrix
      product states}, beginning of Chapter 4 (MPS).
\end{itemize}

\section{Computing entanglement entropy with variational quantum states}
\label{sec:entanglement-entropy-vmc}

TODO: add renyi entropy
\subsection{Direct computation using the SVD}

The entanglement entropy for the ground state can be computed directly as follows:
\begin{enumerate}
\item Obtain an approximation of the ground state
      $\ket{\psi_0}=\sum_\alpha c_\alpha\ket{\alpha}\in\mathcal H^{AB} = \mathcal H^{A} \otimes \mathcal H^{B} $ as some parametrized variational state $\psi_\theta(\alpha)$.
\item Use the variational state (e.g. $\psi_\theta(\alpha) = c_\alpha$) to evaluate the amplitudes on every allowed configuration $\alpha$ and normalize the resulting coefficient vector.
\item Reshape the coefficients into the matrix
      $\psi\in\mathbb C^{d_A\times d_B}$, respecting the ordering of the
      product basis.
\item Compute the SVD of $\psi$ to obtain the Schmidt values $s_i$, yielding the entropy $S=-\sum_i s_i^2\log s_i^2$
\end{enumerate}
This approach is also straightforward with exact diagonalization since the ground state is represented by a vector consisting of the desired amplitudes. However, this quickly becomes unfeasible for ED and variational states as the Hilbert space grows exponentially with the system size.


\subsection{Constructing the photon reduced density matrix}

Given that $d_A\ll d_B$, one can estimate the small reduced density matrix $\rho_A$ and diagonalize it directly. For illustrative purposes, let subsystem $A$ describe a single photon mode and let subsystem $B$ describe a fermionic system.

Let $n=0,\dots,N_{\mathrm{ph}}$ denote the photon occupation and $\sigma$
a fermionic occupation configuration. Define the wavefunction amplitudes by
\begin{equation}
\psi(n,\sigma)\equiv\braket{n,\sigma}{\psi}.
\end{equation}
For a normalized state, the matrix elements of $\rho_A$ follow from \cref{eq:reduced_density_matrix_as_operator,eq:density_matrix_as_wavefunction_coefficients}
\begin{equation}
(\rho_A)_{nn'}
=\sum_\sigma\psi(n,\sigma)\psi^*(n',\sigma)
\end{equation}

Since variational states do not necessarily provide normalized amplitudes, we might have to normalize them in the above formula. Defining
\begin{equation}
Z\equiv\sum_{m,\sigma}|\psi(m,\sigma)|^2,
\end{equation}
we can write the normalized reduced density matrix more generally as
\begin{equation}
\boxed{\,
(\rho_A)_{nn'}
=\frac{\sum_\sigma\psi(n,\sigma)\psi^*(n',\sigma)}{Z}
\,}.
\end{equation}
For a normalized state, $Z=1$.

\subsection{A Monte Carlo estimator for the reduced density matrix}

For each fermionic configuration $\sigma$, define
\begin{equation}
u_n(\sigma)\equiv
\frac{\psi(n,\sigma)}
{\sqrt{\sum_{m=0}^{N_{\mathrm{ph}}}|\psi(m,\sigma)|^2}},
\qquad
\vec u(\sigma)=
\begin{pmatrix}
u_0(\sigma)\\
\vdots\\
u_{N_{\mathrm{ph}}}(\sigma)
\end{pmatrix}
\end{equation}
only if the denominator is non-zero. This is almost always the case when using Monte Carlo sampling since these configurations will never be visited, unless you start in one of them.
This vector is normalized for every $\sigma$ with non-zero marginal
probability:
\begin{equation}
\vec u(\sigma)^\dagger\vec u(\sigma)
=\sum_{n=0}^{N_{\mathrm{ph}}}|u_n(\sigma)|^2=1.
\end{equation}
The matrix elements of its outer product are
\begin{equation}
\bigl(\vec u(\sigma)\vec u(\sigma)^\dagger\bigr)_{nn'}
=
\frac{\psi(n,\sigma)\psi^*(n',\sigma)}
{\sum_{m=0}^{N_{\mathrm{ph}}}|\psi(m,\sigma)|^2}.
\end{equation}

Now consider the marginal probability distribution of the fermionic
configurations,
\begin{align}
p(\sigma) &=
\sum_{m=0}^{N_{\mathrm{ph}}} p(m, \sigma) \\
&\equiv \frac{\sum_{m=0}^{N_{\mathrm{ph}}}|\psi(m,\sigma)|^2}
{\sum_{m,\sigma'}|\psi(m,\sigma')|^2} .
\end{align}
Taking the expectation value with respect to this distribution yields
\begin{align}
\mathbb E_{\sigma\sim p(\sigma)}
\bigl[\vec u(\sigma)\vec u(\sigma)^\dagger\bigr]_{nn'}
&=\sum_\sigma
  \frac{\psi(n,\sigma)\psi^*(n',\sigma)}
       {\sum_m|\psi(m,\sigma)|^2}\,p(\sigma)\notag\\
&=\sum_\sigma
  \frac{\psi(n,\sigma)\psi^*(n',\sigma)}
       {\sum_m|\psi(m,\sigma)|^2}\,
  \frac{\sum_m|\psi(m,\sigma)|^2}
       {\sum_{m,\sigma'}|\psi(m,\sigma')|^2}\notag\\
&=\frac{\sum_\sigma\psi(n,\sigma)\psi^*(n',\sigma)}
       {\sum_{m,\sigma'}|\psi(m,\sigma')|^2}\notag\\
&=(\rho_A)_{nn'}.
\end{align}
Therefore,
\begin{equation}
\boxed{\,
\rho_A=
\mathbb E_{\sigma\sim p(\sigma)}
\bigl[\vec u(\sigma)\vec u(\sigma)^\dagger\bigr]
\,}.
\end{equation}
Moreover, since we are only interested in an expectation value over $\sigma$, we have
\begin{align}
\mathbb E_{(n,\sigma)\sim p(n, \sigma)}
\bigl[F(\sigma)\bigr]
&=\sum_{n,\sigma}
  p(n, \sigma) F(\sigma) \notag\\
&=\sum_{\sigma}
  p(\sigma) F(\sigma) \notag\\
&= \mathbb E_{\sigma\sim p(\sigma)}
\bigl[F(\sigma)\bigr] . \label{eq:expectation_value_equal_probability_marginal_probability}
\end{align}
Thus, we can sample either from $p(\sigma)$ or from $p(n,\sigma)$, depending on what is more convenient.

\subsection{Computation using variational Monte Carlo}

The resulting procedure is:
\begin{enumerate}
\item Sample a configuration $(\tilde n,\sigma)$ using the VMC sampler,
      i.e.\ from the joint distribution
      \begin{equation}
      p(n,\sigma)=\frac{|\psi(n,\sigma)|^2}{Z}.
      \end{equation}
      Discard $\tilde n$ and retain $\sigma$. Note, that $Z$ does not need to be known in VMC. The retained configurations follow the marginal distribution $p(\sigma)$ (see \cref{eq:expectation_value_equal_probability_marginal_probability}).
\item Evaluate all amplitudes $\psi(n,\sigma)$ for
      $n=0,\dots,N_{\mathrm{ph}}$ and construct $\vec u(\sigma)$.
\item Compute $\vec u(\sigma)\vec u(\sigma)^\dagger$.
      Repeat the preceding steps and average the resulting matrices:
      \begin{equation}
      \widehat\rho_A=
      \frac{1}{M}\sum_{k=1}^{M}
      \vec u(\sigma_k)\vec u(\sigma_k)^\dagger
      \approx\rho_A.
      \end{equation}
      The wide hat indicates that we only obtain an estimate of the true density matrix.
\item Diagonalize $\widehat\rho_A$ and use its eigenvalues
      $\widehat\lambda_i$ to estimate the von Neumann entropy:
      \begin{equation}
      \widehat S=-\sum_i\widehat\lambda_i\log\widehat\lambda_i.
      \end{equation}
\end{enumerate}

\subsection{References}
I did not use any literature explicitly, except for AI helping me figure out what to do. The following paper employs a similar approach to the one I demonstrate here.
\begin{itemize}
\item Tang, Yifan, et al. \emph{Deep quantum Monte Carlo approach for polaritonic chemistry.} The Journal of Chemical Physics 163.3 (2025). (Appendix A.1 and A.3)
\end{itemize}

\end{document}
